% Part: first-order-logic
% Chapter: sequent-calculus
% Section: proving-things

\documentclass[../../../include/open-logic-section]{subfiles}

\begin{document}

\iftag{FOL}
      {\olfileid[hi]{fol}{seq}{pro}}
      {\olfileid[hi]{pl}{seq}{pro}}

\olsection{\usetoken{p}{derivation} बनाने के उदाहरण}

\begin{ex}
एक $\Log{LK}$-!!{derivation} दीजिए जिसका अंतिम अनुक्रम
$!A \land !B \Sequent !A$ हो।

आरंभ में वांछित अंतिम अनुक्रम को !!{derivation} के सबसे नीचे लिखते हैं।
\begin{prooftree}
\AxiomC{}
\UnaryInf$!A\land !B \fCenter !A$
\end{prooftree}
अब निर्धारित करना है कि किस प्रकार के अनुमान का निचला अनुक्रम इस रूप का
हो सकता है। यह कोई संरचनात्मक नियम हो सकता है, पर आरंभ किसी तार्किक नियम
को खोजने से करना उचित है। निचले अनुक्रम में आने वाला एकमात्र तार्किक संयोजक
$\land$ है, इसलिए हमें $\land$-नियम चाहिए; और चूँकि $\land$ चिह्न पूर्वपक्ष
में है, हमें \LeftR{\land} नियम देखना है।
\begin{prooftree}
\AxiomC{}
\RightLabel{\LeftR{\land}}
\UnaryInf$!A\land !B \fCenter !A$
\end{prooftree}
\LeftR{\land} अनुमान का ऊपरी अनुक्रम क्या रहा होगा, इसके दो विकल्प हैं:
ऊपरी अनुक्रम $!A
\Sequent !A$ हो सकता है, या $!B \Sequent !A$। स्पष्ट है
कि $!A \Sequent !A$ प्रारंभिक अनुक्रम है (जो हमारे लिए अच्छा है), जबकि
$!B \Sequent !A$ सामान्यतः व्युत्पन्न नहीं है। अब ऊपरी अनुक्रम भर देते हैं:
\begin{prooftree}
\Axiom$!A \fCenter !A$
\RightLabel{\LeftR{\land}}
\UnaryInf$!A\land !B \fCenter !A$
\end{prooftree}
अब हमारे पास सही $\Log{LK}$-!!{derivation} है, जिसका अंतिम अनुक्रम $!A
\land !B \Sequent !A$ है।
\end{ex}

\begin{ex}
एक $\Log{LK}$-!!{derivation} दीजिए जिसका अंतिम अनुक्रम $\lnot !A \lor !B
\Sequent !A \lif !B$ हो।

वांछित अंतिम अनुक्रम को !!{derivation} के सबसे नीचे लिखकर आरंभ कीजिए।
\begin{prooftree}
\AxiomC{}
\UnaryInf$\lnot !A \lor !B \fCenter !A \lif !B$
\end{prooftree}
इस अंतिम अनुक्रम को देने वाला तार्किक नियम खोजने के लिए इसमें आए तार्किक
संयोजक देखते हैं: $\lnot$, $\lor$, और $\lif$। अभी हमारा ध्यान केवल $\lor$
और $\lif$ पर है, क्योंकि अभी हमें !!{main operator}s से काम है और ये अंतिम
अनुक्रम के !!{sentence}s में बाहर से लागू होते हैं; $\lnot$ किसी दूसरे
संयोजक के परास के भीतर है, इसलिए उसे बाद में सँभालेंगे। अतः अंतिम अनुमान के
लिए तार्किक नियमों के विकल्प \LeftR{\lor} नियम और \RightR{\lif} नियम हैं।
वास्तव में दोनों में से कोई भी चुन सकते हैं, पर \RightR{\lif} नियम चुनते हैं
(यदि और कोई कारण न हो, तो इसलिए कि इससे दो शाखाओं में विभाजन कुछ देर टलता
है)। निचला अनुक्रम देने वाले \RightR{\lif} अनुमानों के रूप के अनुसार यह ऐसा
दिखना चाहिए:
\begin{prooftree}
\AxiomC{}
\UnaryInf$ !A, \lnot !A \lor !B \fCenter !B $
\RightLabel{\RightR{\lif}} \UnaryInf$ \lnot !A \lor !B \fCenter !A \lif !B $
\end{prooftree}

यदि पूर्वपक्ष में $\lnot !A \lor !B$ को बाहर की ओर ले आएँ, तो
\LeftR{\lor} नियम लगा सकते हैं। नियम-रूप के अनुसार इसे निम्न प्रकार दो ऊपरी
अनुक्रमों में बँटना चाहिए:
\begin{prooftree}
\AxiomC{}
\UnaryInf$\lnot !A, !A \fCenter !B$
\AxiomC{}
\UnaryInf$!B, !A \fCenter !B$
\RightLabel{\LeftR{\lor}}
\BinaryInf$ \lnot !A \lor !B, !A \fCenter !B $
\RightLabel{\RightR{\Exchange}}
\UnaryInf$ !A, \lnot !A \lor !B \fCenter !B $
\RightLabel{\RightR{\lif}}
\UnaryInf$ \lnot !A \lor !B \fCenter !A \lif !B $
\end{prooftree}
याद रखें कि हम ऊपर की ओर रास्ता बनाते हुए प्रारंभिक अनुक्रमों तक पहुँचने का
प्रयत्न कर रहे हैं; हम लगभग पहुँच गए हैं!{} दाहिनी शाखा प्रारंभिक अनुक्रम से
केवल एक दुर्बलीकरण और एक विनिमय दूर है; इनके बाद वह पूरी हो जाती है:
\begin{prooftree}
\AxiomC{}
\UnaryInf$\lnot !A, !A \fCenter !B$
\Axiom$!B \fCenter !B$
\RightLabel{\LeftR{\Weakening}}
\UnaryInf$!A, !B \fCenter !B$
\RightLabel{\LeftR{\Exchange}}
\UnaryInf$!B, !A \fCenter !B$
\RightLabel{\LeftR{\lor}}
\BinaryInf$\lnot !A \lor !B, !A \fCenter !B $
\RightLabel{\RightR{\Exchange}}
\UnaryInf$ !A, \lnot !A \lor !B \fCenter !B $
\RightLabel{\RightR{\lif}}
\UnaryInf$ \lnot !A \lor !B \fCenter !A \lif !B $
\end{prooftree}

अब बाईं शाखा देखें। यहाँ किसी !!{sentence} में मिलने वाला एकमात्र तार्किक
संयोजक $\lnot$ है, और वह पूर्वपक्ष के !!{sentence}s में आता है; इसलिए हमें
\LeftR{\lnot} नियम का एक प्रयोग चाहिए।
\begin{prooftree}
\AxiomC{}
\UnaryInf$ !A \fCenter !B, !A$
\RightLabel{\LeftR{\lnot}}
\UnaryInf$\lnot !A, !A \fCenter !B$
\Axiom$!B \fCenter !B$
\RightLabel{\LeftR{\Weakening}}
\UnaryInf$!A, !B \fCenter !B$
\RightLabel{\LeftR{\Exchange}}
\UnaryInf$!B, !A \fCenter !B$
\RightLabel{\LeftR{\lor}}
\BinaryInf$\lnot !A \lor !B, !A \fCenter !B $
\RightLabel{\RightR{\Exchange}}
\UnaryInf$ !A, \lnot !A \lor !B \fCenter !B $
\RightLabel{\RightR{\lif}}
\UnaryInf$ \lnot !A \lor !B \fCenter !A \lif !B $
\end{prooftree}
जिस प्रकार दाहिनी शाखा पूरी की थी, उसी प्रकार इस बाईं शाखा को पूरा करने के
लिए भी केवल एक दुर्बलीकरण और एक विनिमय चाहिए।
\begin{prooftree}
\Axiom$!A \fCenter !A$
\RightLabel{\RightR{\Weakening}}
\UnaryInf$ !A \fCenter !A, !B$
\RightLabel{\RightR{\Exchange}}
\UnaryInf$ !A \fCenter !B, !A$
\RightLabel{\LeftR{\lnot}}
\UnaryInf$\lnot !A, !A \fCenter !B$
\Axiom$!B \fCenter !B$
\RightLabel{\LeftR{\Weakening}}
\UnaryInf$!A, !B \fCenter !B$
\RightLabel{\LeftR{\Exchange}}
\UnaryInf$!B, !A \fCenter !B$
\RightLabel{\LeftR{\lor}}
\BinaryInf$\lnot !A \lor !B, !A \fCenter !B $
\RightLabel{\RightR{\Exchange}}
\UnaryInf$ !A, \lnot !A \lor !B \fCenter !B $
\RightLabel{\RightR{\lif}}
\UnaryInf$ \lnot !A \lor !B \fCenter !A \lif !B $
\end{prooftree}
\end{ex}

\begin{ex}
एक $\Log{LK}$-!!{derivation} दीजिए जिसका अंतिम अनुक्रम $\lnot !A \lor \lnot !B
\Sequent \lnot (!A \land !B)$ हो।

ऊपर की विधियों का प्रयोग करते हुए वांछित अंतिम अनुक्रम को सबसे नीचे लिखकर
आरंभ करते हैं।
\begin{prooftree}
\AxiomC{}
\UnaryInf$ \lnot !A \lor \lnot !B \fCenter \lnot (!A \land !B) $
\end{prooftree}
अंतिम अनुक्रम के !!{sentence}s में उपलब्ध प्रधान संक्रियक $\lor$ चिह्न और
$\lnot$ चिह्न हैं। यहाँ \LeftR{\lor} या \RightR{\lnot}, दोनों में से कोई
नियम लगाया जा सकता है; पर हम \RightR{\lnot} से आरंभ करते हैं, क्योंकि इससे
कुछ देर तक दो शाखाओं में बँटना टल जाता है:
\begin{prooftree}
\AxiomC{}
\UnaryInf$!A \land !B, \lnot !A \lor \lnot !B \fCenter $
\RightLabel{\RightR{\lnot}}
\UnaryInf$\lnot !A \lor \lnot !B \fCenter \lnot (!A \land !B)$
\end{prooftree}
अब \LeftR{\land} और \LeftR{\lor} में से किस नियम को देखें, यह चुनना है।
पहले देखें कि \LeftR{\land} नियम लगाने पर क्या होता है: ऊपरी अनुक्रम के
लिए $!A,
\lnot !A \lor !B \Sequent \quad$ और $!B, \lnot !A
\lor !B \Sequent \quad$ में से कोई एक चुन सकते हैं। चूँकि !!{derivation}, $!A$ और
$!B$ के संबंध में सममित है, पहला विकल्प लेते हैं:
\begin{prooftree}
\AxiomC{}
\UnaryInf$!A, \lnot !A \lor \lnot !B \fCenter $
\RightLabel{\LeftR{\land}}
\UnaryInf$!A \land !B, \lnot !A \lor \lnot !B \fCenter $
\RightLabel{\RightR{\lnot}}
\UnaryInf$\lnot !A \lor \lnot !B \fCenter \lnot (!A \land !B)$
\end{prooftree}
!!{derivation} को आगे भरने पर एक समस्या सामने आती है:
\begin{prooftree}
\Axiom$!A \fCenter !A$
\RightLabel{\LeftR{\lnot}}
\UnaryInf$ \lnot !A, !A \fCenter$
\AxiomC{}
\RightLabel{?}
\UnaryInf$!A \fCenter !B$
\RightLabel{\LeftR{\lnot}}
\UnaryInf$ \lnot !B, !A \fCenter$
\RightLabel{\LeftR{\lor}}
\BinaryInf$\lnot !A \lor \lnot !B, !A \fCenter $
\RightLabel{\LeftR{\Exchange}}
\UnaryInf$!A, \lnot !A \lor \lnot !B \fCenter $
\RightLabel{\LeftR{\land}}
\UnaryInf$!A \land !B, \lnot !A \lor \lnot !B \fCenter $
\RightLabel{\RightR{\lnot}}
\UnaryInf$\lnot !A \lor \lnot !B \fCenter \lnot (!A \land !B)$
\end{prooftree}
दाहिनी शाखा के शीर्ष को और घटाया नहीं जा सकता, और संरचनात्मक अनुमानों द्वारा
उसे प्रारंभिक अनुक्रम तक भी नहीं लाया जा सकता; इसलिए यह सही रास्ता नहीं है।
अर्थात् ऊपर \LeftR{\land} नियम लगाना भूल थी। पहले वाली स्थिति पर लौटकर उसके
बदले \LeftR{\lor} नियम लगाने पर मिलता है:
\begin{prooftree}
\AxiomC{}
\UnaryInf$\lnot !A, !A \land !B \fCenter $

\AxiomC{}
\UnaryInf$\lnot !B, !A \land !B \fCenter $

\RightLabel{\LeftR{\lor}}
\BinaryInf$\lnot !A \lor \lnot !B, !A \land !B \fCenter $
\RightLabel{\LeftR{\Exchange}}
\UnaryInf$!A \land !B, \lnot !A \lor \lnot !B \fCenter $
\RightLabel{\RightR{\lnot}}
\UnaryInf$\lnot !A \lor \lnot !B \fCenter \lnot (!A \land !B)$
\end{prooftree}
हर शाखा को पहले की तरह पूरा करने पर मिलता है:
\begin{prooftree}
\Axiom$ !A \fCenter!A$
\RightLabel{\LeftR{\land}}
\UnaryInf$!A \land !B \fCenter !A$
\RightLabel{\LeftR{\lnot}}
\UnaryInf$\lnot !A, !A \land !B \fCenter $

\Axiom$ !B \fCenter !B$
\RightLabel{\LeftR{\land}}
\UnaryInf$!A \land !B \fCenter !B$
\RightLabel{\LeftR{\lnot}}
\UnaryInf$\lnot !B, !A \land !B \fCenter $

\RightLabel{\LeftR{\lor}}
\BinaryInf$\lnot !A \lor \lnot !B, !A \land !B \fCenter $
\RightLabel{\LeftR{\Exchange}}
\UnaryInf$!A \land !B, \lnot !A \lor \lnot !B \fCenter $
\RightLabel{\RightR{\lnot}}
\UnaryInf$\lnot !A \lor \lnot !B \fCenter \lnot (!A \land !B)$
\end{prooftree}
(इन चरणों में $\land$ नियमों को $\lnot$ नियमों से नीचे भी लगा सकते थे और तब
भी सही !!{derivation} मिलती।)
\end{ex}

\begin{ex}
अब तक हमने संकुचन नियम का प्रयोग नहीं किया, पर कभी-कभी वह आवश्यक होता है।
यह उसका एक उदाहरण है। मान लें कि हमें $\quad \Sequent !A \lor \lnot !A$
सिद्ध करना है। $\RightR{\lor}$ को उल्टी दिशा में लगाने पर इन दो
!!{derivation}s में से कोई एक मिलेगी:
\begin{prooftree}
\AxiomC{}
\UnaryInf$ \fCenter !A$
\RightLabel{\RightR{\lor}}
\UnaryInf$ \fCenter !A \lor \lnot !A$
\DisplayProof\qquad\bottomAlignProof
\AxiomC{}
\UnaryInf$!A \fCenter $
\RightLabel{\RightR{\lnot}}
\UnaryInf$ \fCenter \lnot !A$
\RightLabel{\RightR{\lor}}
\UnaryInf$ \fCenter !A \lor \lnot !A$
\end{prooftree}
निस्संदेह इनमें से कोई भी प्रारंभिक अनुक्रम पर समाप्त नहीं होती। युक्ति यह
समझना है कि संकुचन नियम !!a{sentence} की दो प्रतियों को एक में मिला देता
है---और प्रमाण खोजते समय, अर्थात् नीचे से ऊपर जाते हुए, पूर्वाधार में $!A \lor \lnot
!A$ की एक प्रति रख सकते हैं; जैसे:
\begin{prooftree}
\AxiomC{}
\UnaryInf$ \fCenter !A \lor \lnot !A, !A$
\RightLabel{\RightR{\lor}}
\UnaryInf$ \fCenter !A \lor \lnot !A, !A \lor \lnot !A$
\RightLabel{\RightR{\Contraction}}
\UnaryInf$ \fCenter !A \lor \lnot !A$
\end{prooftree}
अब $\RightR{\lor}$ को दूसरी बार लगा सकते हैं और साथ ही~$\lnot
!A$ पा सकते हैं; इससे पूरी !!{derivation} मिल जाती है।
\begin{prooftree}
\Axiom$!A \fCenter !A$
\RightLabel{\RightR{\lnot}}
\UnaryInf$\fCenter !A, \lnot !A$
\RightLabel{\RightR{\lor}}
\UnaryInf$\fCenter !A, !A \lor \lnot !A$
\RightLabel{\RightR{\Exchange}}
\UnaryInf$ \fCenter !A \lor \lnot !A, !A$
\RightLabel{\RightR{\lor}}
\UnaryInf$ \fCenter !A \lor \lnot !A, !A \lor \lnot !A$
\RightLabel{\RightR{\Contraction}}
\UnaryInf$ \fCenter !A \lor \lnot !A$
\end{prooftree}
\end{ex}

\begin{prob}
निम्नलिखित अनुक्रमों की !!{derivation}s दीजिए:
\begin{enumerate}
\item $!A \land (!B \land !C) \Sequent (!A \land !B) \land !C$.
\item $!A \lor (!B \lor !C) \Sequent (!A \lor !B) \lor !C$.
\item $!A \lif (!B \lif !C) \Sequent !B \lif (!A \lif !C)$.
\item $!A \Sequent \lnot\lnot !A$.
\end{enumerate}
\end{prob}

\begin{prob}
निम्नलिखित अनुक्रमों की !!{derivation}s दीजिए:
\begin{enumerate}
\item $(!A \lor !B) \lif !C \Sequent !A \lif !C$.
\item $(!A \lif !C) \land (!B \lif !C) \Sequent (!A \lor !B) \lif !C$.
\item $\Sequent \lnot(!A \land \lnot !A)$.
\item $!B \lif !A \Sequent \lnot !A \lif \lnot !B$.
\item $\Sequent (!A \lif \lnot !A) \lif \lnot !A$.
\item $\Sequent \lnot(!A \lif !B) \lif \lnot !B$.
\item $!A \lif !C \Sequent \lnot (!A \land \lnot !C)$.
\item $!A \land \lnot !C \Sequent \lnot (!A \lif !C)$.
\item $!A \lor !B, \lnot !B \Sequent !A$.
\item $\lnot !A \lor \lnot !B \Sequent \lnot(!A \land !B)$.
\item $\Sequent (\lnot !A \land \lnot !B) \lif\lnot(!A \lor !B)$.
\item $\Sequent \lnot(!A \lor !B) \lif (\lnot !A \land \lnot !B)$.
\end{enumerate}
\end{prob}

\begin{prob}
निम्नलिखित अनुक्रमों की !!{derivation}s दीजिए:
\begin{enumerate}
\item $\lnot(!A \lif !B) \Sequent !A$.
\item $\lnot(!A \land !B) \Sequent \lnot !A \lor \lnot !B$.
\item $!A \lif !B \Sequent \lnot !A \lor !B$.
\item $\Sequent \lnot \lnot !A \lif !A$.
\item $!A \lif !B, \lnot !A \lif !B \Sequent !B$.
\item $(!A \land !B) \lif !C \Sequent (!A \lif !C) \lor (!B \lif !C)$.
\item $(!A \lif !B) \lif !A \Sequent !A$.
\item $\Sequent (!A \lif !B) \lor (!B \lif !C)$.
\end{enumerate}
(इन सभी में $\RightR{\Contraction}$~नियम आवश्यक है।)
\end{prob}
\end{document}
